\section{Mục tiêu 2: NTK Regime}
\begin{frame}{Mục tiêu 2: NTK regime khi chiều rộng $m \to \infty$}
  Từ Mục tiêu 1, ta có:
  \[
    \frac{d\mathbf f(t)}{dt}
    = -K_m(\theta_m(t))(\mathbf f(t) - \mathbf y),
  \]
  trong đó $K_m$ là NTK thực nghiệm của mạng rộng $m$.

  \vspace{0.2cm}
  \textbf{Vấn đề:} $K_m(\theta_m(t))$ vẫn phụ thuộc vào $t$
  vì tham số $\theta_m(t)$ thay đổi khi train.
\end{frame}

\begin{frame}{Mục tiêu 2:}
      \textbf{Mục tiêu 2:} Chứng minh rằng trong giới hạn $m \to \infty$,
  \begin{itemize}
    \item Tại $t=0$, $K_m(\theta_m(0))$ hội tụ về một kernel giới hạn
          \[
            \Theta(X,X) \quad (\text{deterministic}).
          \]
    \item Trong suốt quá trình train trên khoảng thời gian hữu hạn $t \in [0,T]$,
          $K_m(\theta_m(t))$ gần như \textbf{không đổi}:
          \[
            K_m(\theta_m(t)) \approx K_m(\theta_m(0)).
          \]
  \end{itemize}

  \vspace{0.2cm}
  Gộp lại:
  \[
    K_m(\theta_m(t))
    \xrightarrow[m\to\infty]{}
    \Theta(X,X),
  \]
  và gần như \textbf{không phụ thuộc $t$} (NTK regime / lazy training).
\end{frame}

\begin{frame}{$g_{ir}(0)$ có bậc $O(1/\sqrt{m})$}
  Nhắc lại mô hình:
  \[
    f_m(x;\theta)
    = \frac{1}{\sqrt{m}} \sum_{r=1}^m a_r\,\sigma(w_r^\top x).
  \]

  Xét neuron thứ $r$, tham số:
  \[
    \theta_r = (a_r, w_r).
  \]
  Gradient của output tại điểm $x_i$ theo neuron $r$:
  \[
    g_{ir}(0) := \nabla_{\theta_r} f_m(x_i; \theta(0)).
  \]
\end{frame}

\begin{frame}{$g_{ir}(0)$ có bậc $O(1/\sqrt{m})$}
      \textbf{Đạo hàm theo $a_r$:}
  \[
    \frac{\partial f_m(x_i)}{\partial a_r}
    = \frac{1}{\sqrt{m}} \sigma(w_r^\top x_i).
  \]
Vậy:
  \[
    \frac{\partial f_m}{\partial a_r}
    = O\!\left(\frac{1}{\sqrt{m}}\right).
  \]

  \vspace{0.2cm}
  \textbf{Đạo hàm theo $w_r$:}
  \[
    \frac{\partial f_m(x_i)}{\partial w_r}
    = \frac{1}{\sqrt{m}}\,a_r\,\sigma'(w_r^\top x_i)\, x_i.
  \]
Vậy:
  \[
    \left\|\frac{\partial f_m}{\partial w_r}\right\|
    = O\!\left(\frac{1}{\sqrt{m}}\right).
  \]
\end{frame}

\begin{frame}{$g_{ir}(0)$ có bậc $O(1/\sqrt{m})$}
  \textbf{Kết luận:}
  \[
    g_{ir}(0)
    =
    \Big(
      \frac{\partial f_m}{\partial a_r},
      \frac{\partial f_m}{\partial w_r}
    \Big),
    \quad
    \Rightarrow\quad
    \|g_{ir}(0)\|
    = O\!\left(\frac{1}{\sqrt{m}}\right).
  \]

  \textbf{Hệ quả quan trọng:}
  \begin{itemize}
    \item Mỗi neuron $r$ đóng góp một lượng nhỏ vào gradient khi $m$ lớn.
    \item Đóng góp vào NTK của neuron $r$ tại cặp $(x_i,x_j)$ là
          \[
            \big\langle g_{ir}(0), g_{jr}(0) \big\rangle
            = O\!\left(\frac{1}{m}\right),
          \]
    \item Tổng qua $m$ neuron:
          \[
            \sum_{r=1}^m O\!\left(\frac{1}{m}\right)
            = O(1)
          \]
  \end{itemize}
\end{frame}

\begin{frame}{Mục tiêu 2 – Bước 1: hội tụ của NTK tại $t=0$}
  Jacobian tại init:
  \[
    J_m(0) := \nabla_\theta \mathbf f_m(\theta_m(0)).
  \]
NTK tại init:
  \[
    K_m(\theta_m(0)) = J_m(0)J_m(0)^\top.
  \]

  Mỗi phần tử:
  \[
    K_{m,ij}(0)
      = \sum_{r=1}^m \big\langle g_{ir}(0), g_{jr}(0)\big\rangle,
  \]
  với $g_{ir}(0) := \nabla_{\theta_r} f_m(x_i;\theta_m(0))$ như đã phân tích.
\end{frame}

\begin{frame}{Mục tiêu 2 – Bước 1: luật số lớn cho NTK}
  Do các neuron $r$ là i.i.d.,
  \[
    K_{m,ij}(0)
    = \sum_{r=1}^m \big\langle g_{ir}(0),g_{jr}(0)\big\rangle
    \xrightarrow[m\to\infty]{a.s.} \Theta(x_i,x_j).
  \]

  \textbf{Kết luận Bước 1:}
  \[
    K_m(\theta_m(0)) \xrightarrow[m\to\infty]{a.s.} \Theta(X,X).
  \]
  Tức là ngay tại khởi tạo, NTK thực nghiệm hội tụ về một kernel không phụ thuộc vào m.
\end{frame}

\begin{frame}{Mục tiêu 2 – Bước 2: tham số thay đổi rất ít khi train}
  \textbf{Ý tưởng của Bước 2:}  
  Khi $m$ rất lớn, mỗi tham số chỉ thay đổi rất nhỏ quanh điểm khởi tạo
  $\theta_m(0)$ trong một khoảng thời gian hữu hạn $[0,T]$.

  \vspace{0.2cm}
  Ta dùng gradient flow:
  \[
    \frac{d\theta_r}{dt}
    = -\nabla_{\theta_r}L(\theta),
    \quad r = 1,\dots,m.
  \]
\end{frame}

\begin{frame}{Mục tiêu 2 – Bước 2:}
  \begin{itemize}
    \item Với mỗi $r$, $L(\theta) = \tfrac12\sum_i (f_m(x_i)-y_i)^2$.
\item Ta có:
\[
\begin{aligned}
\frac{\partial L}{\partial a_r}
&= \sum_{i=1}^n (f_m(x_i)-y_i)
      \frac{\partial f_m(x_i)}{\partial a_r} \\
&= \sum_{i=1}^n (f_i-y_i)\frac{1}{\sqrt{m}}\sigma(w_r^\top x_i)
 = O\!\left(\frac{1}{\sqrt{m}}\right).
\end{aligned}
\]

    \item Tương tự,
      \[
        \frac{\partial L}{\partial w_r}
        = \sum_{i=1}^n (f_i-y_i)
            \frac{1}{\sqrt{m}} a_r \sigma'(w_r^\top x_i)x_i
        = O\!\left(\frac{1}{\sqrt{m}}\right).
      \]
  \end{itemize}

  \[ \Rightarrow
    \frac{d\theta_r}{dt}
    = -\nabla_{\theta_r}L(\theta)
    = O\!\left(\frac{1}{\sqrt{m}}\right),
    \quad \forall r.
  \]
\end{frame}

\begin{frame}{Mục tiêu 2 – Bước 2: kernel ``đóng băng''}
  Với $t \in [0,T]$ hữu hạn:
  \[
    \|\theta_r(t) - \theta_r(0)\|
    \le \int_0^T \left\| \frac{d\theta_r}{ds} \right\| ds
    = O\!\left(\frac{T}{\sqrt{m}}\right)
    \xrightarrow[m\to\infty]{} 0.
  \]

  \vspace{0.2cm}
  Bây giờ xét NTK trong quá trình train:
  \[
    K_m(\theta(t)) = J_m(\theta(t))J_m(\theta(t))^\top.
  \]
  Xét phần tử $(i,j)$:
  \[
    K_{m,ij}(t)
      = \sum_{r=1}^m \big\langle g_{ir}(t), g_{jr}(t)\big\rangle,
    \quad
    g_{ir}(t) := \nabla_{\theta_r} f_m(x_i;\theta(t)).
  \]
\end{frame}

\begin{frame}{Khái niệm Lipschitz và vai trò trong NTK}
\textbf{Định nghĩa (Lipschitz):}
Một hàm $h(\theta)$ được gọi là \emph{Lipschitz} nếu tồn tại hằng số $L$ sao cho
\[
  \|h(\theta_1) - h(\theta_2)\|
  \;\le\;
  L \, \|\theta_1 - \theta_2\|,
  \quad \forall\theta_1,\theta_2.
\]
$L$ gọi là \textbf{Lipschitz constant}.

\vspace{0.3cm}
\textbf{Ý nghĩa:}
\begin{itemize}
  \item Hàm không được “nhảy đột ngột”; thay đổi trong đầu vào kéo theo thay đổi có kiểm soát ở đầu ra.
  \item Nếu đạo hàm của $h$ bị chặn bởi $M$ (tức $|\nabla h|\le M$), thì $h$ Lipschitz với hằng số $L\le M$.
\end{itemize}
\end{frame}

\begin{frame}{Áp dụng trong NTK}
\textbf{Áp dụng vào NTK:}
\[
g_{ir}(\theta_r)
= \nabla_{\theta_r} f_m(x_i;\theta)
\]
có dạng
\[
g_{ir} =
\frac{1}{\sqrt{m}}
\begin{pmatrix}
\sigma(w_r^\top x_i) \\
a_r \sigma'(w_r^\top x_i) x_i
\end{pmatrix}.
\]

Nhờ $\sigma$ trơn và $\sigma',\sigma''$ bị chặn:
\[
\|\nabla_{\theta_r} g_{ir}\| = O\!\left(\frac{1}{\sqrt{m}}\right)
\quad\Rightarrow\quad
g_{ir} \text{ là Lipschitz theo }\theta_r.
\]

Do đó:
\[
\|g_{ir}(t) - g_{ir}(0)\|
\le
O\!\left(\frac{1}{\sqrt{m}}\right)
\cdot
\|\theta_r(t)-\theta_r(0)\|
=
O\!\left(\frac{1}{m}\right).
\]
\end{frame}

\begin{frame}{Bước 2: thay đổi của một hạng tử NTK}
Đặt:
\[
  k_{ij}^{(r)}(t) := \langle g_{ir}(t),\, g_{jr}(t)\rangle.
\]

Hiệu giữa hai thời điểm:
\[
  |k_{ij}^{(r)}(t) - k_{ij}^{(r)}(0)|
  \le
  \|g_{ir}(t)-g_{ir}(0)\|\cdot \|g_{jr}(t)\|
  \;+\;
  \|g_{jr}(t)-g_{jr}(0)\|\cdot \|g_{ir}(0)\|.
\]

Do:
\[
  \|g_{ir}(t)-g_{ir}(0)\| = O(1/m), \qquad
  \|g_{jr}(t)\| = O(1/\sqrt{m}),
\]

ta được:
\[
  |k_{ij}^{(r)}(t) - k_{ij}^{(r)}(0)|
  = O\!\left(\frac{1}{m}\right)
    \cdot O\!\left(\frac{1}{\sqrt{m}}\right)
  = O\!\left(\frac{1}{m\sqrt{m}}\right).
\]
\end{frame}

\begin{frame}{Bước 2: tổng hợp tác động của toàn bộ $m$ neuron}
Tổng thay đổi của NTK:
\[
  |K_{m,ij}(t) - K_{m,ij}(0)|
  = \left|\sum_{r=1}^m \big(k_{ij}^{(r)}(t) - k_{ij}^{(r)}(0)\big)\right|.
\]

Mỗi hạng tử có bậc $O(1/(m\sqrt{m}))$, nên:
\[
  |K_{m,ij}(t) - K_{m,ij}(0)|
  \le
  m \cdot O\!\left(\frac{1}{m\sqrt{m}}\right)
  = O\!\left(\frac{1}{\sqrt{m}}\right)
  \xrightarrow[m\to\infty]{}
  0.
\]

\textbf{Kết luận Bước 2:}
\[
  \|K_m(\theta_m(t)) - K_m(\theta_m(0))\|
  = O\!\left(\frac{1}{\sqrt{m}}\right)
  \to 0.
\]
\end{frame}

\begin{frame}{Mục tiêu 2 – Kết luận NTK regime}
  Gộp Bước 1 và Bước 2:
  \begin{itemize}
    \item (Bước 1) Tại khởi tạo,
      \[
        K_m(\theta_m(0)) \xrightarrow[m\to\infty]{a.s.} \Theta(X,X).
      \]
    \item (Bước 2) Trong quá trình train với thời gian hữu hạn,
      \[
        \sup_{t\in[0,T]}
        \|K_m(\theta_m(t)) - K_m(\theta_m(0))\|
        \xrightarrow[m\to\infty]{} 0.
      \]
  \end{itemize}

  \vspace{0.2cm}
  \textbf{Kết luận:} Trong NTK regime, ta có xấp xỉ:
  \[
    \frac{d\mathbf f(t)}{dt}
      = -K_m(\theta_m(t))(\mathbf f(t)-\mathbf y)
      \approx -\Theta(X,X)(\mathbf f(t)-\mathbf y).
  \]
\end{frame}
